% sections/introduction.tex
\chapter{Introduzione}
\label{chap:introduction}

\section{Il sistema AIS}\label{sec:intro_ais}

L'Automatic Identification System (AIS) è un sistema di trasmissione e ricezione dati utilizzato in maniera
diffusa sia da imbarcazioni commerciali che da diporto.
Nasce con lo scopo di migliorare la sicurezza della navigazione riducendo il rischio di collisioni e,
con l'evolversi della tecnologia, assume anche un ruolo legato al tracciamento delle imbarcazioni.
Dal 2004 l'International Maritime Organization (IMO) ne ha reso obbligatoria l'installazione su tutte le navi
passeggeri e sulle navi commerciali al di sopra di una certa stazza~\cite{imo2004}.

Al di là di queste finalità operative, l'AIS rappresenta oggi una fonte ricca di informazioni per lo studio
dei comportamenti delle navi. Ogni nave trasmette periodicamente la propria posizione, la velocità,
la direzione del moto, lo stato di navigazione e altre caratteristiche statiche, come tipo e dimensioni.
La frequenza di trasmissione varia con la velocità e lo stato della nave, da pochi secondi a qualche minuto;
i dati utilizzati in questa tesi, come si vedrà, sono invece campionati a un record al minuto.
Grazie alla larga adozione del sistema, anche aree circoscritte ma trafficate generano grandi moli di dati,
che per vastità, rumorosità e complessità richiedono tecniche automatiche per essere sfruttate.

\section{Il problema dei comportamenti anomali}\label{sec:intro_problem}

Tra le informazioni che si possono estrarre dai dati AIS, una delle più utili è la segnalazione dei
comportamenti anomali, cioè dei movimenti che si discostano da quelli abituali.
Una nave che abbandona il corridoio percorso di solito tra due porti, che si ferma in un punto insolito
o che conclude il proprio viaggio lontano da un porto può essere in difficoltà, oppure coinvolta in attività
illecite come il contrabbando o la pesca in aree protette.
Questa tesi si concentra sulla prima di queste situazioni, la \emph{route deviation} (deviazione dalla rotta),
che è anche la tipologia di anomalia più studiata in letteratura; le altre categorie sono descritte nel
Capitolo~\ref{chap:literature}.

Riconoscere queste situazioni a mano non è praticabile: un solo trimestre di dati relativi a un arcipelago
contiene milioni di messaggi, e in generale non esiste un elenco di casi già classificati come normali o anomali
da cui partire.
Servono quindi metodi \emph{non supervisionati}, in grado di apprendere dai soli dati storici che cosa sia
il comportamento normale e di segnalare ciò che se ne discosta.

\section{L'approccio: il framework TREAD}\label{sec:intro_approach}

Il metodo adottato è il framework TREAD (Traffic Route Extraction and Anomaly Detection), proposto
da Pallotta et al.~\cite{pallotta} e pubblicato nel 2013.
L'idea alla base è semplice. Dai dati storici si individuano dapprima i \emph{waypoint}, cioè le aree in cui
le navi si fermano oppure entrano o escono dalla zona osservata, come porti e ancoraggi.
I percorsi seguiti dalle navi da un waypoint all'altro definiscono poi le \emph{rotte}, e ogni singolo viaggio
compiuto da una nave lungo una rotta prende il nome di \emph{trip}.
A ogni trip viene infine assegnato un punteggio che risponde a due domande: i punti del trip si trovano dove
passano di solito le navi di quella rotta? Il moto da un punto al successivo è prevedibile a partire dalla
posizione e dalla velocità precedenti?
Un trip che passa per zone poco frequentate, o che si muove in modo difficile da prevedere, ottiene un punteggio
che lo distingue dai trip normali.

Il metodo risale a più di dieci anni fa ed è stato sviluppato con le tecnologie dell'epoca.
L'obiettivo della tesi è riproporlo, re-implementandolo con gli strumenti oggi disponibili e applicandolo
a un dataset AIS di grandi dimensioni. Il metodo è ripreso nella sua struttura generale, mentre alcuni passaggi
sono stati adattati o migliorati.

\section{Obiettivi e quesiti di ricerca}\label{sec:intro_questions}

I quesiti di ricerca su cui è stato sviluppato il lavoro sono i seguenti:
\begin{enumerate}
    \item Il framework TREAD, sviluppato agli inizi degli anni 2010, è ancora applicabile con gli strumenti attuali
          a un dataset AIS reale di milioni di record?
    \item Quali aspetti e caratteristiche dell'algoritmo sono ``riciclabili'' e quali hanno richiesto adattamenti o sostituzioni?
    \item Quali passaggi del metodo possono essere migliorati con gli strumenti attuali, e con quali effetti?
    \item Quali metriche è opportuno utilizzare per quantificare le prestazioni dell'algoritmo?
\end{enumerate}

\section{Caso di studio e contributi}\label{sec:intro_case}

I dati utilizzati sono i messaggi AIS pubblicati dal NOAA (National Oceanic and Atmospheric Administration),
l'agenzia statunitense che tra i suoi compiti ha il monitoraggio oceanico.
Sono stati preferiti ad altre fonti perché sottoposti dall'ente a un controllo di qualità, e perché gli attributi
statici delle navi sono verificati sul registro ufficiale della Guardia Costiera statunitense~\cite{noaa_ais2019}.
L'area di studio è l'arcipelago delle Hawaii, una zona circoscritta ma con un traffico intenso ed eterogeneo,
che va dai pescherecci alle navi mercantili e passeggeri, fino ai rimorchiatori e alle imbarcazioni da diporto.
Il periodo analizzato copre i mesi da gennaio a marzo 2019. La durata di tre mesi è dettata dall'hardware disponibile;
l'anno 2019 è stato scelto in quanto ultimo prima della pandemia di Covid-19, così da evitare possibili effetti
della pandemia sul traffico marittimo.
Il metodo viene infine valutato su una singola rotta, quella tra il porto di Honolulu e il porto di Kalaeloa,
entrambi sull'isola di Oahu. Tra le rotte più frequentate è stata scelta questa perché percorsa in modo molto
regolare, condizione che rende ben definito il comportamento normale e quindi più riconoscibili le deviazioni.

Rispetto al lavoro originale, la tesi apporta i seguenti contributi:
\begin{itemize}
    \item la re-implementazione dell'intera pipeline in Python, con librerie moderne per la gestione di grandi
          volumi di dati;
    \item l'identificazione dei waypoint tramite H-DBSCAN, in grado di gestire aree a densità di traffico molto
          diversa, al posto di DBSCAN;
    \item un calcolo diretto delle componenti della velocità, che evita di gestire l'ambiguità di quadrante
          dell'arcotangente;
    \item una valutazione quantitativa del punteggio tramite un insieme di trip etichettati manualmente,
          assente nel lavoro originale.
\end{itemize}

La tesi è stata svolta nell'ambito del progetto europeo KOIOS\footnote{\url{https://www.koiosproject.eu/}},
i cui fondi hanno finanziato la proposta di questo argomento di tesi.

\section{Struttura della tesi}\label{sec:intro_structure}

Il resto della tesi è organizzato come segue.
\begin{itemize}
    \item Il Capitolo~\ref{chap:literature} descrive il funzionamento del sistema AIS, le tipologie di anomalia
          e i principali approcci di anomaly detection presenti in letteratura, tra cui il framework TREAD.
    \item Il Capitolo~\ref{chap:methodology} descrive il dataset, il pre-processing dei dati, l'implementazione
          delle fasi di identificazione dei waypoint e delle rotte e la fase di anomaly detection, evidenziando
          le differenze rispetto al lavoro originale.
    \item Il Capitolo~\ref{chap:results} presenta i risultati ottenuti sul dataset delle Hawaii e la valutazione
          del punteggio di anomalia sulla rotta selezionata.
    \item Il Capitolo~\ref{chap:conclusions} riassume il lavoro svolto, ne discute i limiti e indica i possibili
          sviluppi futuri.
\end{itemize}
